\section{Method}
\label{sec:method}

\begin{figure*}[!t]
\centering
\includegraphics[width=0.98\linewidth]{figures/fig2_topoprm_framework.pdf}
\caption{Overview of our proposed \method with its three-stage framework \framework. The DAG Scorer lifts reasoning traces into structural representations. A hierarchical multiplicative reward and asymmetric advantage clipping bound within-stratum credit, while a topology-conditioned on-policy distillation stage compresses the long-CoT teacher into a compact student.}
\label{fig:overview}
\end{figure*}

\framework (Figure~\ref{fig:overview}) is a topology-aware post-training framework whose three stages share a \textbf{DAG Scorer} $\mathcal{E}$ that converts long-chain reasoning traces into a structure-sensitive form on which all downstream supervision is computed. Stage~I performs a supervised cold start to enforce parsable formatting; Stage~II runs GRPO with a hierarchical reward that fuses \textbf{TopoPRM} with Stratified Clipping Advantage Estimation (\estimation); Stage~III executes Topology-Guided Self-Distillation, in which the Stage-II teacher provides token-level reverse-KL supervision while the student rolls out under topology-conditioned revision prompts.

\subsection{Problem Setup}
Let $x=(q,c)$ denote an input pairing a mathematical problem $q$ with optional context $c$, and let $\pi_\theta$ be a policy that generates a reasoning trace $y=(s_1,\dots,s_T)\sim\pi_\theta(\cdot\!\mid\!x)$, where each $s_i$ is a textual reasoning step. Outcome-only RLVR supervises $y$ with a sparse final-answer reward $r_{\mathrm{out}}(x,y)\!\in\![0,1]$. We define $\mathcal{G}_y=(\mathcal{V},\mathcal{E}_d)$ as a directed dependency graph on $y$, with node set $\mathcal{V}=\{s_1,\dots,s_T\}$ and edge set $\mathcal{E}_d\subseteq\mathcal{V}\!\times\!\mathcal{V}$ encoding surface-evidenced support relations between steps, and $\mathcal{E}:y\!\mapsto\!\mathcal{G}_y$ as the DAG Scorer that produces it. \textit{Our objective is to recover latent dependency structure from $\mathcal{G}_y$ as unified supervision across SFT, RL, and distillation, without human step labels or learned verifier models.}

\subsection{Topology-Aware Data Processing}
\label{sec:reward_module}

\paragraph{Trace-to-DAG Extraction.}
Given a rollout $y=(s_1,\dots,s_T)$, the DAG Scorer $\mathcal{E}$ proceeds in three coupled stages whose outputs are reused by every downstream signal. \textbf{(i)~Step segmentation and parsing.} $\mathcal{E}$ first splits $y$ into reasoning steps using numbered markers, mathematical discourse markers, and line boundaries, then partitions the trace into sub-questions so that no edge can cross independent components. \textbf{(ii)~Implicit feature extraction.} Each step $s_i$ is reduced to three normalized sets, mathematical expressions $E_i$, canonicalised claim keys $C_i$, and variable references $V_i$, after information-density filters discard single symbols, pure numerals, and heading-like fragments that would otherwise inflate the matching surface. \textbf{(iii)~DAG construction.} For every ordered intra-component pair $(s_i,s_j)$ with $i\!<\!j$, $\mathcal{E}$ inserts an implicit edge $i\!\to\!j$ whenever $E_i\!\cap\!E_j$, $C_i\!\cap\!C_j$, or $V_i\!\cap\!V_j$ is non-empty; if $s_j$ is a derivation or conclusion step still lacking incoming support, an \emph{adaptive sequential weak} edge from $s_{j-1}$ is added only when $s_{j-1}$ does not already emit strong outgoing dependencies, which prevents the chain-dominant artefacts that long traces tend to produce. Full extraction rules are deferred to Appendix~\ref{app:dag}.

\paragraph{Topology and Continuity Diagnostics.}
\label{sec:diagnostics}
On the resulting graph $\mathcal{G}_y\!=\!(\mathcal{V},\mathcal{E}_d)$ we compute two normalized scores that constitute the process-level reward signal. The \emph{topology quality} $q_{\mathrm{topo}}\!\in\![0,1]$ aggregates five structural checks, while the \emph{continuity quality} $q_{\mathrm{cont}}\!\in\![0,1]$ measures step-level traceability:
\begin{align}
q_{\mathrm{topo}}(\mathcal{G}_y) &= \mathbb{1}[|\mathcal{V}|>0]
+\mathbb{1}[\mathcal{G}_y\text{ acyclic}] \notag\\
+\mathbb{1}[\rho_{\mathrm{orph}}&(\mathcal{G}_y)\!=\!0]+\delta(\mathcal{G}_y)
+\kappa(\mathcal{G}_y,\mathcal{G}^{\star}), \label{eq:qtopo}\\[2pt]
q_{\mathrm{cont}}(y) &= \gamma\,\eta(y)+(1-\gamma)\,\mathbb{1}[\eta(y)\!=\!1], \label{eq:qcont}\\[2pt]
r_{\mathrm{topo}}(y) &= \alpha\,q_{\mathrm{topo}}(\mathcal{G}_y)+(1-\alpha)\,q_{\mathrm{cont}}(y), \label{eq:rtopo}
\end{align}
where $\rho_{\mathrm{orph}}$, $\delta$, and $\kappa$ measure, respectively, the fraction of orphan conclusion nodes, the proportion of edges respecting the premise-to-consequence ordering, and the alignment with an optional reference DAG $\mathcal{G}^{\star}$; $\eta(y)=\tfrac{1}{T}\sum_{i=1}^{T}\mathbb{1}[\mathrm{supported}(s_i,s_{<i},x)]$
is the fraction of reasoning steps in $y$ whose entities overlap with the prefix $s_{<i}$ or the question $x$, and $\gamma,\alpha\!\in\!(0,1)$ are fixed mixing coefficients reported in Appendix~\ref{app:hparam}.

\paragraph{Hierarchical Aggregation.}
Weighted-sum rewards can mis-rank traces when high process scores compensate for incorrect final answers~\citep{Denison2024SycophancyTS}, and they tend to collapse the within-group variance on which GRPO relies~\citep{yu2025dapo,kazemnejad2025vineppo}. We therefore couple the correctness-based and structural signals multiplicatively, treating $r_{\mathrm{topo}}$ as a structural modulator on top of an outcome-anchored base:
\begin{equation}
\label{eq:r_total}
\begin{aligned}
r_{\mathrm{total}}(x,y)=\;
&\bigl(w_o\,r_{\mathrm{out}}+w_f\,r_{\mathrm{fmt}}+w_l\,r_{\mathrm{len}}\bigr)\\
&\cdot\bigl(1+r_{\mathrm{topo}}(y)\bigr),
\end{aligned}
\end{equation}
where $r_{\mathrm{out}}\!\in\![0,1]$ measures final-answer correctness, $r_{\mathrm{fmt}}$ scores parsability, and $r_{\mathrm{len}}$ regularises length. The modulator lies in $[1,2]$, so the multiplicative form preserves correctness primacy: a trace with $r_{\mathrm{out}}\!=\!0$ cannot be lifted by any process bonus, while traces with equal outcome scores remain distinguishable through their dependency structure. This yields the within-group reward dispersion that \estimation exploits in Stage~II.

\subsection{Topology-Guided Post-Training}
\label{sec:post_training}
% The three stages of \framework share a single DAG-based supervision interface.
% Stage~I (SFT) aligns format and basic reasoning behaviour, Stage~II (GRPO with hierarchical reward) injects topology-sensitive process supervision without overriding answer correctness, and Stage~III (\framework) reuses the same diagnostics for compact-student distillation.

\paragraph{Stage I: Supervised Warm-Start.}
We perform Supervised Fine-Tuning (SFT) on structured mathematical reasoning data to serve as a warm-start that teaches the model to follow the required response format, produce stepwise solutions, and expose intermediate reasoning in a form parsable by $\mathcal{E}$. This reduces invalid rollouts during subsequent reinforcement learning and stabilizes reward computation, since topology diagnostics are meaningful only when the generated trace contains recoverable reasoning steps.

\paragraph{Stage II: GRPO with TopoPRM Rewards.}
Starting from the SFT checkpoint, we optimise the policy with GRPO under the TopoPRM reward. For each prompt, the policy samples a group of completions, each of which is parsed into a dependency DAG by $\mathcal{E}$ and scored accordingly. The aggregated reward in Eq.~\ref{eq:r_total} combines outcome correctness, format compliance, length regularization, and topology-derived process signals. Outcome rewards select for correct final answers, while topology and continuity rewards encourage the model to make intermediate dependencies explicit, avoid unsupported conclusions, and reduce structurally redundant reasoning.
% This stage is where TopoPRM directly improves reasoning behavior. Outcome rewards select for correct final answers, while topology and continuity rewards encourage the model to make intermediate dependencies explicit, avoid unsupported conclusions, and reduce structurally redundant reasoning. SCAE further constrains the auxiliary structural rewards so that they improve credit assignment within correctness strata rather than rewarding plausible-looking but wrong traces.

\paragraph{Anti-collapse mechanisms.}
Multiplicative structural rewards introduce a potential hacking risk, because an incorrect response may appear well structured under surface-form diagnostics. Following \citet{liu2026graph}, we adopt Stratified Clipping Advantage Estimation (\estimation), which partitions each GRPO group by correctness and bounds the auxiliary contribution of $q_{\mathrm{topo}}$ and $q_{\mathrm{cont}}$ within each stratum, so that structural scores cannot flip the sign of an outcome-wrong trace:
\begin{equation}
\label{eq:scae}
\widehat{A}_i \;=\;
\mathrm{clip}\!\Big(\tfrac{R_i-\mu_{\mathcal{B}^{\pm}}}{\sigma_{\mathcal{B}^{\pm}}+\epsilon},\,
\,\underline{c}^{\pm},\,\overline{c}^{\pm}\Big),
\quad i\in\mathcal{B}^{\pm}.
\end{equation}
We further mitigate reward collapse with batch-relative normalization, correctness-first shaping, and a positive base floor for low-variance groups, so that process gains do not vanish when outcome rewards are sparse. Implementation details and ablations are reported in Appendix~\ref{app:scae}.

% To reduce reward collapse in GRPO, we use batch-relative normalisation and correctness-first shaping, and monitor low-variance groups while retaining a positive base floor so that process gains do not vanish when outcome rewards are sparse. As an additional safeguard against multiplicative reward hacking, we follow \citet{liu2026graph} and estimate advantages with stratified clipping (\estimation), which partitions each GRPO group by correctness and bounds the auxiliary contribution of $q_{\mathrm{topo}}$ and $q_{\mathrm{cont}}$ within each stratum so that structural scores cannot flip the sign of an outcome-wrong trace:
% \begin{equation}
% \label{eq:scae}
% \widehat{A}_i \;=\;
% \mathrm{clip}\!\Big(\tfrac{R_i-\mu_{\mathcal{B}^{\pm}}}{\sigma_{\mathcal{B}^{\pm}}+\epsilon},\,
% \,\underline{c}^{\pm},\,\overline{c}^{\pm}\Big),
% \qquad i\in\mathcal{B}^{\pm}.
% \end{equation}
% Implementation and ablation of \estimation, together with diagnostics for the anti-collapse monitors, are reported in Appendix~\ref{app:scae}.

% \paragraph{Advantage estimation with SCAE.}
% Structural rewards introduce a potential reward-hacking risk: an incorrect response may appear well structured under surface-form diagnostics. We therefore adopt Stratified Clipping Advantage Estimation (SCAE) from prior topology-aware RL work and adapt it to our reward setting. For each GRPO group, completions are partitioned into correctness strata according to \(r_{\mathrm{out}}\). Within each stratum, auxiliary topology and continuity rewards may re-rank samples, but clipping prevents them from flipping the sign of an incorrect response's advantage. Thus, SCAE acts as a safeguard that keeps structural credit assignment subordinate to final-answer correctness.

\paragraph{Stage~III: Topology-Guided Self Distillation.}
\label{sec:tgsd}
The optimised model after Stage~II serves as the teacher $\pi_T$, while a smaller pretrained checkpoint initialises the student $\pi_\theta$. \framework performs on-policy distillation conditioned on topology diagnostics. For each prompt $x$, the student first rolls out its own response $y^{\star}\!\sim\!\pi_\theta(\cdot\!\mid\!x)$; the same extractor $\mathcal{E}$ produces $\mathcal{G}_{y^{\star}}$ together with $(r_{\mathrm{out}},q_{\mathrm{topo}},q_{\mathrm{cont}})$; a deterministic dispatcher $\psi$ then issues a revision instruction $P_r=\psi(r_{\mathrm{out}},q_{\mathrm{topo}},k)$ targeting an orphan-conclusion position $k$ in $\mathcal{G}_{y^{\star}}$. The teacher then supervises the student under $P_r$, so that the signal concentrates on structurally suspect regions (orphan conclusions and continuity breaks) rather than spreading uniformly across the trace. Sampling $y$ from $\pi_\theta$ rather than from a fixed teacher trace also eliminates the train--test mismatch of off-policy distillation~\citep{agarwal2024gkd,song2026opdsurvey}.

\paragraph{Distillation objective.}
Let $p_t \!\triangleq\! \pi_\theta(\cdot\!\mid\!x,y_{<t})$ and $q_t \!\triangleq\! \pi_T(\cdot\!\mid\!x,y^{\star},P_r,y_{<t})$ denote the student and teacher next-token distributions at position $t$, where $y_{<t}$ is the autoregressive prefix sampled on-policy by the student. Restricted to samples passing format and topology filters ($r_{\mathrm{out}}^{(\mathrm{rev})}\!=\!1$ and $q_{\mathrm{topo}}^{(\mathrm{rev})}\!\ge\!\tau_g$), the student minimises
\begin{equation}
\label{eq:tgsd}
\mathcal{L}_{\mathrm{TGSD}}(\theta)=
\mathbb{E}_{x,\,y\sim\pi_\theta}\!\sum_{t=1}^{|y|}
D_{\mathrm{KL}}\!\bigl(p_t\,\|\,q_t\bigr),
\end{equation}
where $\mathcal{G}_{y^{\star}}$ exposes broken support, $P_r$ steers $\pi_T$ toward the missing dependency, turning the per-token KL into a structurally targeted gradient.

\paragraph{Critical dependency chain as the compression target.}
Equation~\eqref{eq:tgsd} pushes the student toward a linearisation of $\mathcal{G}_{y^{\star}}$ obtained by composing three operations on the DAG: \emph{step-type contraction} merges consecutive same-type nodes; \emph{layer aggregation} folds parallel topological generations into a layered chain $L_1\!\to\!L_2\!\to\!\cdots$; and \emph{transitive sparsification} replaces $\mathcal{G}$ by its transitive reduction $\mathrm{tr}(\mathcal{G})$. The resulting \emph{critical dependency chain}, analogous to the critical path in project scheduling, lower-bounds any structurally supported trajectory that closes the answer; the student is therefore trained to recover this minimal chain rather than to match teacher token frequencies, which is the structural counterpart of length-penalty compression. The full \method pipeline is summarised in Algorithm~\ref{alg:topoprm}.

\begin{algorithm}[t]
\small
\caption{\framework Training Pipeline}
\label{alg:topoprm}
\begin{algorithmic}[1]
\Require Base policy $\pi_0$; data $\mathcal{D}_{\mathrm{sft}},\mathcal{D}_{\mathrm{rl}}$; weights $w_o,w_f,w_l,\alpha$
\Ensure Teacher $\pi_T$, compact student $\pi_S$
\State $\pi_\theta\!\leftarrow\!\mathrm{SFT}(\pi_0,\mathcal{D}_{\mathrm{sft}})$ \Comment{Stage~I}
\For{each GRPO step on $\mathcal{D}_{\mathrm{rl}}$} \Comment{Stage~II}
  \State Sample $\{y^{(i)}\}_{i=1}^{G}\!\sim\!\pi_\theta(\cdot\!\mid\!x)$; build $\mathcal{G}^{(i)}\!=\!\mathcal{E}(y^{(i)})$
  \State Compute $r_{\mathrm{out}}^{(i)},r_{\mathrm{fmt}}^{(i)},r_{\mathrm{len}}^{(i)},q_{\mathrm{topo}}^{(i)},q_{\mathrm{cont}}^{(i)}$
  \State $r_{\mathrm{total}}^{(i)}$ by Eq.~\ref{eq:r_total}; $\widehat{A}^{(i)}$ by Eq.~\ref{eq:scae}
  \State Update $\pi_\theta$ with the GRPO objective
\EndFor
\State $\pi_T\!\leftarrow\!\pi_\theta$; init.\ $\pi_S$ from a smaller pretrained checkpoint
\For{each prompt $x\!\in\!\mathcal{D}_{\mathrm{rl}}$} \Comment{Stage~III}
  \State $y_{\mathrm{init}}\!\sim\!\pi_S(\cdot\!\mid\!x)$; compute $(r_{\mathrm{out}},q_{\mathrm{topo}},q_{\mathrm{cont}})$ on $\mathcal{E}(y_{\mathrm{init}})$
  \State $P_r\!\leftarrow\!\psi(r_{\mathrm{out}},q_{\mathrm{topo}},k)$
  \State $y_{\mathrm{rev}}\!\sim\!\pi_T(\cdot\!\mid\!x,y_{\mathrm{init}},P_r)$; filter by format/topology
  \State Update $\pi_S$ with $\mathcal{L}_{\mathrm{TGSD}}$ (Eq.~\ref{eq:tgsd})
\EndFor
\State \Return $(\pi_T,\pi_S)$
\end{algorithmic}
\end{algorithm}